% TOP_PROBLEM: {"id": 9700022, "title": "Percolation of planar empires at unit merger rate", "category_id": 19, "category": {"id": 19, "name": "probability", "display_name": "Probability", "description": "Problems involving probability theory, stochastic processes, and random structures.", "slug": "probability", "order_index": 19, "created_at": "2026-07-31T15:26:25.671Z"}, "status": "open", "problem_number": "AMR-096-0022", "set_id": 13}
% FIRST_POSTED: 2026-10-01
% ENTRY_KIND: statement_only
% UnsolvedMath ID 9700022; source catalogue code AMR-096-0022
% !TeX program = lualatex
% Definition and sources only; this file is not an LLM solution attempt.
\documentclass[11pt,a4paper]{article}
\usepackage[margin=25mm]{geometry}
\usepackage{fontspec}
\setmainfont{lmroman10-regular.otf}[BoldFont=lmroman10-bold.otf,
 ItalicFont=lmroman10-italic.otf,BoldItalicFont=lmroman10-bolditalic.otf]
\usepackage{amsmath,amssymb,mathtools}
\usepackage{unicode-math}
\setmathfont{latinmodern-math.otf}
\AtBeginDocument{\let\setminus\smallsetminus}
\usepackage{xurl,hyperref}
\hypersetup{colorlinks=true,urlcolor=blue}
\setcounter{secnumdepth}{0}
\setlength{\parindent}{0pt}
\setlength{\parskip}{5pt}
\setlength{\emergencystretch}{3em}
\begin{document}
\section*{Percolation of planar empires at unit merger rate}\label{9700022}
{\small UnsolvedMath ID 9700022 · AMR-096-0022 · Probability · AMR Open Problem Lists}
\subsection{Definitions and mathematical statement}
Start with a locally finite partition of $\mathbb R^2$ into polygonal cells, using one of the source's initial models: a regular square, triangular, or hexagonal tessellation, or the Voronoi tessellation of a homogeneous planar Poisson point process. In the last case the points form a random locally finite set with independent Poisson counts in disjoint bounded regions.

Call current regions empires. Two empires are adjacent if their boundaries share a segment of positive length, rather than only a point. In continuous time, every unordered adjacent pair merges into its union at rate one. Equivalently, conditional on the current configuration, an adjacent pair has probability $dt+o(dt)$ of merging during a short interval of length $dt$; after mergers, rates are assigned to the newly adjacent pairs.

Does this process percolate: is there a finite time at which an unbounded empire exists? Determine the answer for the indicated initial tessellations, in particular the deterministic square grid.

The rate is per pair of empires, independently of their areas, perimeters, and the number of original cell edges along their common boundary. Thus many shared boundary segments do not multiply that pair's merger rate. Empires can become nonconvex and may acquire holes; connectedness is preserved by each merger.

The target concerns formation at finite time of a region containing infinitely many original cells. It is not enough that bounded empires grow without bound as time tends to infinity. In an infinite-volume construction, possible emergence of such a region is precisely the percolation event under study.
\subsection{Short English statement}
Do planar regions merging at rate one per adjacent pair form an unbounded empire in finite time?
\subsection{Sources}
\begin{itemize}
\item[S1] ULAM.AI and the UnsolvedMath Contributors, supplied problems.json, record 9700022 (AMR-096-0022), AMR Open Problem Lists. Catalogue identity and source transcription; CC BY 4.0.
\url{https://huggingface.co/datasets/ulamai/UnsolvedMath}.
\item[S2] Aldous - Open problems, empires.html, second Problem. Original source indexed by AMR-096-0022.
\url{https://www.stat.berkeley.edu/~aldous/Research/OP/index.html}.
\item[S3] D. Aldous, Empires and Percolation, original rate-one question and initial configurations.
\url{https://www.stat.berkeley.edu/~aldous/Research/OP/empires.html}.
\end{itemize}
\end{document}
